\chapter{The Mealy Comonad, Behaviour, and Nerode Realization}
\label{ch:mealy-comonad-behaviour-nerode}

\section{Overview}
\label{sec:mealy-comonad-overview}

The preceding chapters developed internal categories as monads in
\(\mathbf{Span}(\mathcal E)\), Mealy morphisms as monad morphisms in spans,
and relative program categories as action categories of restricted
Eilenberg--Moore actions. This chapter develops the corresponding
coalgebraic and behavioural theory.

The main point is that Mealy morphisms are not merely coalgebras for a raw
transition-output endofunctor. They are coalgebras for a comonad naturally
constructed from the forgetful functor from Mealy systems to their carrier
spans.

Fix internal categories
\[
\mathbb A
=
(A_0,A_1,s_A,t_A,e_A,m_A)
\qquad\text{and}\qquad
\mathbb B
=
(B_0,B_1,s_B,t_B,e_B,m_B)
\]
in the sense of Chapter~\ref{ch:spans-internal-categories-mealy}. We keep the
same composition convention: if
\[
a:u\to v,
\qquad
b:v\to w,
\]
then
\[
b\circ a:u\to w
\]
means first \(a\), then \(b\). Equivalently,
\[
m_A(a,b)=b\circ a.
\]

The behavioural development proceeds in layers.

First, under a locally presentable locally cartesian closed hypothesis, we
construct a right adjoint to the forgetful functor
\[
U_{\mathbb A,\mathbb B}:
\mathsf{Mealy}(\mathbb A,\mathbb B)
\longrightarrow
\mathcal E/(A_0\times B_0).
\]
This gives the Mealy comonad
\[
\mathbb M_{\mathbb A,\mathbb B}
:=
U_{\mathbb A,\mathbb B}R_{\mathbb A,\mathbb B}.
\]
Its Eilenberg--Moore coalgebras are canonically equivalent to Mealy morphisms
\[
\mathbb A\rightsquigarrow\mathbb B.
\]

\[
\begin{tikzcd}[column sep=large]
\mathsf{Mealy}(\mathbb A,\mathbb B)
\arrow[r, shift left=1.1ex, "U_{\mathbb A,\mathbb B}"]
&
\mathcal E/(A_0\times B_0)
\arrow[l, shift left=1.1ex, "R_{\mathbb A,\mathbb B}"]
\end{tikzcd}
\qquad
\mathbb M_{\mathbb A,\mathbb B}=U_{\mathbb A,\mathbb B}R_{\mathbb A,\mathbb B}.
\]

Second, the final behaviour system is constructed as the cofree Mealy system
on the terminal carrier:
\[
\mathsf{Beh}_{\mathbb A,\mathbb B}
:=
R_{\mathbb A,\mathbb B}(1_{A_0\times B_0}).
\]
Thus finality is not assumed. It is a consequence of the cofree construction.

The cofree construction is then unpacked into a residual form. A behaviour
over an input object \(v\in A_0\) and output object \(y\in B_0\) is a coherent
assignment of output interfaces and output arrows to all residual input
histories ending at \(v\). In generalized-element notation, this is a pointed
residual-slice functor
\[
\mathbb A/v\longrightarrow\mathbb B
\]
whose value at the identity arrow \(1_v:v\to v\) is \(y\). This residual-slice
description is the internal-categorical analogue of the usual final-behaviour
description of deterministic Mealy machines.

In the external one-object case, where
\[
\mathcal E=\mathbf{Set},
\qquad
\mathbb A=M,
\qquad
\mathbb B=N
\]
are monoids, the residual-slice description becomes the set of
\(N\)-valued residual cocycles
\[
\lambda:M\times M\to N
\]
satisfying
\[
\lambda(r,1)=1,
\qquad
\lambda(r,mn)=\lambda(r,m)\lambda(rm,n).
\]
This is the native behaviour object for Mealy machines whose inputs are
already monoid elements.

If the input monoid is further specialized to a free monoid
\[
M=I^\ast,
\]
then restriction of such a coherent residual cocycle to generator inputs gives
the classical final Mealy behaviour object
\[
N^{I^+},
\]
where \(I^+\) is the set of nonempty words. This recovers the standard
coalgebraic form of the final coalgebra of the classical deterministic Mealy
functor
\[
X\mapsto (N\times X)^I
\]
\cite{RuttenUniversalCoalgebra,SilvaCoalgebraicAutomata,PattinsonCoalgebraNotes}.

Third, after adding regularity, we construct realized behaviour objects as
images of behaviour maps. After adding Barr-exactness, we identify derivative
closures with Nerode quotients. This gives a general Myhill--Nerode theorem
without any finiteness assumption:
\[
\text{minimal realization}
=
\text{derivative closure}
=
\text{Nerode quotient}.
\]
The finite-state Myhill--Nerode theorem is then a corollary obtained by
choosing a stable class of finite carrier objects.

\section{Ambient hypotheses}
\label{sec:mealy-comonad-ambient-hypotheses}

For the construction of the Mealy comonad, we assume throughout the first part
of this chapter that
\[
\mathcal E
\]
is locally presentable and locally cartesian closed.

This hypothesis supplies two kinds of structure.

First, local cartesian closure provides dependent products. Hence for every
map
\[
f:X\to Y
\]
the pullback functor
\[
f^\ast:\mathcal E/Y\to\mathcal E/X
\]
has a right adjoint
\[
\Pi_f:\mathcal E/X\to\mathcal E/Y.
\]

\[
\begin{tikzcd}[column sep=large]
\mathcal E/Y
\arrow[r, shift left=1.1ex, "f^\ast"]
&
\mathcal E/X
\arrow[l, shift left=1.1ex, "\Pi_f"]
\end{tikzcd}
\]

Second, local presentability gives the adjoint-functor-theorem setting needed
to construct cofree Mealy systems from the forgetful functor.

Later sections add further hypotheses only when needed:
\[
\begin{array}{c|c}
\text{additional structure} & \text{use} \\
\hline
\text{regularity} & \text{images and realized behaviour objects} \\
\text{Barr-exactness} & \text{effective Nerode quotients} \\
\text{stable finite-state class} & \text{finite-state corollary}
\end{array}
\]

\begin{remark}
\label{rem:not-assuming-final-coalgebra}
We do not assume in advance that a final Mealy coalgebra exists. The final
behaviour object is constructed as a cofree Mealy system on the terminal
carrier. The locally presentable locally cartesian closed hypothesis is used
to construct the relevant cofree functor.
\end{remark}

\section{The input-action comonad}
\label{sec:input-action-comonad}

Let
\[
\mathbb A=(A_0,A_1,s_A,t_A,e_A,m_A)
\]
be an internal category.

As in Chapter~\ref{ch:relative-mealy-program-categories}, \(\mathbb A\)
induces the right-action monad
\[
T_{\mathbb A}:\mathcal E/A_0\to\mathcal E/A_0
\]
given by
\[
T_{\mathbb A}(X\xrightarrow{p}A_0)
=
\left(
A_1\times_{t_A,A_0,p}X
\xrightarrow{s_A\pi_1}
A_0
\right).
\]
Equivalently,
\[
T_{\mathbb A}
=
\Sigma_{s_A}t_A^\ast.
\]

Since \(\mathcal E\) is locally cartesian closed, \(T_{\mathbb A}\) has a
right adjoint
\[
G_{\mathbb A}
:=
\Pi_{t_A}s_A^\ast.
\]

\[
\begin{tikzcd}[column sep=large]
\mathcal E/A_0
\arrow[r, shift left=1.1ex, "{T_{\mathbb A}=\Sigma_{s_A}t_A^\ast}"]
&
\mathcal E/A_0
\arrow[l, shift left=1.1ex, "{G_{\mathbb A}=\Pi_{t_A}s_A^\ast}"]
\end{tikzcd}
\]

The monad structure of \(T_{\mathbb A}\), induced by
\[
e_A
\qquad\text{and}\qquad
m_A,
\]
mates across the adjunction
\[
T_{\mathbb A}\dashv G_{\mathbb A}
\]
to a comonad structure on \(G_{\mathbb A}\).

\begin{proposition}
\label{prop:input-action-comonad-coalgebras}
Coalgebras for the comonad \(G_{\mathbb A}\) are precisely right
\(\mathbb A\)-actions.
\end{proposition}

\begin{proof}
Let
\[
p:X\to A_0
\]
be an object of \(\mathcal E/A_0\). A morphism
\[
\gamma:X\to G_{\mathbb A}X
\]
over \(A_0\) transposes along
\[
T_{\mathbb A}\dashv G_{\mathbb A}
\]
to a map
\[
\alpha:
A_1\times_{t_A,A_0,p}X\to X
\]
over \(A_0\). The condition that \(\alpha\) lies over \(A_0\) is
\[
p\alpha=s_A\pi_1.
\]

The counit law for the \(G_{\mathbb A}\)-coalgebra is the mate of the unit
law for the monad \(T_{\mathbb A}\), hence gives
\[
\alpha(e_A(px),x)=x.
\]
The comultiplication law for the coalgebra is the mate of the multiplication
law for \(T_{\mathbb A}\), hence gives
\[
\alpha(m_A(a,b),x)=\alpha(a,\alpha(b,x))
\]
for composable arrows
\[
a:u\to v,
\qquad
b:v\to w,
\qquad
p(x)=w.
\]
These are exactly the right-action laws.
\end{proof}

\begin{example}[The one-object input case]
\label{ex:input-action-comonad-one-object}
If \(\mathcal E=\mathbf{Set}\) and \(\mathbb A=M\) is a one-object internal
category, then \(M\) is a monoid. We use the convention inherited from
Chapter~\ref{ch:spans-internal-categories-mealy} that
\[
mn=m\circ n,
\]
so \(mn\) means first \(n\), then \(m\), as categorical composition.

Because the Mealy input action is a right action, it processes composite
arrows contravariantly relative to categorical composition. Thus the
operational execution of \(mn=m\circ n\) is: process \(m\) first, then process
\(n\). A \(G_M\)-coalgebra is therefore a set \(X\) with a right \(M\)-action
\[
X\times M\to X,
\qquad
(x,m)\mapsto x\cdot m,
\]
satisfying
\[
x\cdot 1=x,
\qquad
x\cdot(mn)=(x\cdot m)\cdot n.
\]
Thus the input comonad remembers the monoid structure of \(M\), not merely the
underlying set of inputs.
\end{example}

\section{Mealy systems as labelled input coalgebras}
\label{sec:mealy-systems-labelled-input-coalgebras}

Fix internal categories
\[
\mathbb A=(A_0,A_1,s_A,t_A,e_A,m_A),
\qquad
\mathbb B=(B_0,B_1,s_B,t_B,e_B,m_B).
\]
Write
\[
Z:=A_0\times B_0.
\]

Let
\[
\mathsf{Mealy}(\mathbb A,\mathbb B)
\]
denote the external category whose objects are the Mealy morphisms
\[
\mathbb A\rightsquigarrow\mathbb B
\]
defined in Chapter~\ref{ch:spans-internal-categories-mealy}, and whose
morphisms are Mealy transformations.

Equivalently, an object of
\[
\mathsf{Mealy}(\mathbb A,\mathbb B)
\]
is a carrier
\[
A_0\xleftarrow{p}P\xrightarrow{q}B_0
\]
together with maps
\[
\delta:
A_1\times_{t_A,A_0,p}P\to P,
\qquad
\omega:
A_1\times_{t_A,A_0,p}P\to B_1,
\]
satisfying the typing equations
\[
p\delta=s_A\pi_1,
\qquad
s_B\omega=q\delta,
\qquad
t_B\omega=q\pi_P,
\]
and the unit and multiplication equations
\[
\delta(e_A(px),x)=x,
\qquad
\omega(e_A(px),x)=e_B(qx),
\]
\[
\delta(m_A(a,b),x)=\delta(a,\delta(b,x)),
\]
\[
\omega(m_A(a,b),x)
=
m_B\bigl(\omega(a,\delta(b,x)),\omega(b,x)\bigr).
\]
Equivalently, since \(m_A(a,b)=b\circ a\), this last equation says
\[
\omega(b\circ a,x)
=
\omega(b,x)\circ\omega(a,\delta(b,x)).
\]

The action component \(\delta\) makes
\[
p:P\to A_0
\]
a right \(\mathbb A\)-action. Its action category
\[
\int_{\mathbb A}P
\]
has object object \(P\), arrow object
\[
A_1\times_{t_A,A_0,p}P,
\]
source and target
\[
s_{\int}(a,x)=x,
\qquad
t_{\int}(a,x)=\delta(a,x).
\]
Thus an arrow \((a,x)\) of \(\int_{\mathbb A}P\) goes from \(x\) to
\(\delta(a,x)\).

The output map \(\omega\) determines an internal functor
\[
\Omega_P:\int_{\mathbb A}P\to\mathbb B^{\mathrm{op}}
\]
with object map
\[
q:P\to B_0.
\]
The arrow map is \(\omega\), regarded as an arrow of
\(\mathbb B^{\mathrm{op}}\). Explicitly, the Mealy typing equations say that
\[
\omega(a,x):q(\delta(a,x))\to q(x)
\]
is an arrow of \(\mathbb B\), hence an arrow
\[
q(x)\to q(\delta(a,x))
\]
of \(\mathbb B^{\mathrm{op}}\). Therefore the source-target equations for
\(\Omega_P\) into \(\mathbb B^{\mathrm{op}}\) are, in terms of the structure
maps of \(\mathbb B\),
\[
t_B\omega=q\pi_P,
\qquad
s_B\omega=q\delta.
\]
The identity and composition preservation equations for \(\Omega_P\) are
exactly the Mealy output unit and multiplication laws.

\begin{definition}[\(\mathbb B^{\mathrm{op}}\)-labelled input coalgebra]
\label{def:Bop-labelled-input-coalgebra}
A \(\mathbb B^{\mathrm{op}}\)-labelled \(G_{\mathbb A}\)-coalgebra is a
right \(\mathbb A\)-action
\[
(P\xrightarrow{p}A_0,\delta)
\]
together with an internal functor
\[
\Omega_P:\int_{\mathbb A}P\to\mathbb B^{\mathrm{op}}.
\]

A morphism
\[
(P,\delta,\Omega_P)\longrightarrow(Q,\delta',\Omega_Q)
\]
of \(\mathbb B^{\mathrm{op}}\)-labelled \(G_{\mathbb A}\)-coalgebras is a
morphism
\[
\theta:P\to Q
\]
over \(A_0\) preserving the right \(\mathbb A\)-actions and satisfying
\[
\Omega_Q\circ \int_{\mathbb A}\theta=\Omega_P.
\]
Equivalently, it is a carrier map preserving both transition and output.
\end{definition}

\begin{proposition}
\label{prop:mealy-labelled-input-coalgebras}
The category
\[
\mathsf{Mealy}(\mathbb A,\mathbb B)
\]
is equivalently the category of \(\mathbb B^{\mathrm{op}}\)-labelled
\(G_{\mathbb A}\)-coalgebras.
\end{proposition}

\begin{proof}
By Proposition~\ref{prop:input-action-comonad-coalgebras}, a
\(G_{\mathbb A}\)-coalgebra is precisely a right \(\mathbb A\)-action. Given
such an action on \(p:P\to A_0\), the associated action category
\[
\int_{\mathbb A}P
\]
has object object \(P\) and arrow object
\[
A_1\times_{t_A,A_0,p}P.
\]

A labelling of this coalgebra by \(\mathbb B^{\mathrm{op}}\) is an internal
functor
\[
\Omega_P:\int_{\mathbb A}P\to\mathbb B^{\mathrm{op}}.
\]
Its object map is a morphism
\[
q:P\to B_0,
\]
and its arrow map is a morphism
\[
\omega:A_1\times_{t_A,A_0,p}P\to B_1.
\]
Because \(\Omega_P\) lands in \(\mathbb B^{\mathrm{op}}\), the functorial
source-target equations are, in terms of \(\mathbb B\),
\[
t_B\omega=q\pi_P,
\qquad
s_B\omega=q\delta.
\]
These are precisely the Mealy typing equations.

Identity and composition preservation by \(\Omega_P\) are precisely the Mealy
unit and multiplication equations for \(\omega\). Morphisms between labelled
coalgebras are precisely carrier maps preserving the right action and the
\(\mathbb B^{\mathrm{op}}\)-labelling, hence precisely Mealy transformations.
\end{proof}

\begin{example}[The one-object Mealy case]
\label{ex:one-object-labelled-coalgebras}
Let \(\mathcal E=\mathbf{Set}\), let \(\mathbb A=M\), and let
\(\mathbb B=N\), where \(M\) and \(N\) are monoids. With the multiplication
convention
\[
mn=m\circ n,
\]
a Mealy system is a right \(M\)-set \(P\) together with a function
\[
\omega:M\times P\to N
\]
satisfying
\[
\omega(1,x)=1
\]
and
\[
\omega(mn,x)=\omega(m,x)\,\omega(n,x\cdot m).
\]
Here multiplication in \(N\) is written with the same convention as
categorical composition in the one-object category \(N\).
\end{example}

\section{The forgetful functor}
\label{sec:mealy-forgetful-functor}

Define
\[
U_{\mathbb A,\mathbb B}:
\mathsf{Mealy}(\mathbb A,\mathbb B)
\longrightarrow
\mathcal E/Z
\]
by sending a Mealy system carried by
\[
A_0\xleftarrow{p}P\xrightarrow{q}B_0
\]
to the object
\[
(p,q):P\to A_0\times B_0.
\]
On morphisms it sends a Mealy transformation to its underlying map of carrier
spans.

\begin{proposition}
\label{prop:mealy-forgetful-creates-colimits}
The functor
\[
U_{\mathbb A,\mathbb B}
\]
creates colimits.
\end{proposition}

\begin{proof}
Let
\[
D:J\to \mathsf{Mealy}(\mathbb A,\mathbb B)
\]
be a diagram, and write its carriers as
\[
P_j\to Z.
\]
Let
\[
P:=\operatorname{colim}_{j\in J}P_j
\]
be the colimit in \(\mathcal E/Z\).

Since \(\mathcal E\) is locally presentable, all small colimits exist in the
relevant slices. Since \(\mathcal E\) is locally cartesian closed, every
pullback functor between slices has a right adjoint. Thus pullback functors
between slices preserve colimits. Consequently,
\[
A_1\times_{t_A,A_0,p}P
\cong
\operatorname{colim}_{j\in J}
\left(
A_1\times_{t_A,A_0,p_j}P_j
\right).
\]

The same argument applies to the finite iterated pullback domains appearing
in the unit and multiplication laws, for example the domains indexed by
composable input pairs
\[
(A_1\times_{s_A,A_0,t_A}A_1)\times_{t_A\pi_2,A_0,p}P.
\]
All these domains are obtained from \(P\) by pullback along fixed maps between
slices, hence also preserve the colimit cocone.

The transition maps of the \(P_j\)'s therefore induce a unique transition map
\[
\delta:
A_1\times_{t_A,A_0,p}P\to P.
\]
Similarly, the output maps induce a unique output map
\[
\omega:
A_1\times_{t_A,A_0,p}P\to B_1.
\]
The typing equations and the Mealy unit and multiplication equations hold
because they hold after precomposition with the colimit cocone maps, and maps
out of a colimit are determined by the corresponding compatible cocones.

It remains to check the universal property. Let \(Q\) be a Mealy system and
let
\[
P_j\to Q
\]
be a cocone of Mealy transformations. Since \(P\) is the colimit of the
underlying carriers in \(\mathcal E/Z\), there is a unique map
\[
P\to Q
\]
in \(\mathcal E/Z\). To see that it is a Mealy transformation, it suffices to
check transition and output preservation after precomposition with the
colimit cocone. Since
\[
A_1\times_{A_0}P_j\to A_1\times_{A_0}P
\]
also forms a colimit cocone, these equations follow from the corresponding
equations for each \(P_j\to Q\).

Thus \(U_{\mathbb A,\mathbb B}\) creates colimits.
\end{proof}

\begin{corollary}
\label{cor:mealy-forgetful-preserves-colimits}
The functor
\[
U_{\mathbb A,\mathbb B}
\]
preserves colimits.
\end{corollary}

\section{Local presentability and cofree Mealy systems}
\label{sec:cofree-mealy-systems}

\begin{lemma}[Accessible dependent-polynomial structures]
\label{lem:accessible-dependent-polynomial-structures}
Let \(\mathcal E\) be locally presentable and locally cartesian closed. Let
\[
\mathcal C
\]
be a locally presentable slice of \(\mathcal E\). Any category of structures
over \(\mathcal C\) specified by accessible dependent-polynomial operations
and finite-limit equations is locally presentable, and its forgetful functor
to \(\mathcal C\) is accessible.
\end{lemma}

\begin{proof}
Dependent-polynomial operations are built from pullback, dependent sum, and
dependent product functors between locally presentable slices. These functors
are accessible under the present hypotheses. Categories obtained by adding
structure maps for accessible operations are constructed from inserters and
comma categories of accessible functors between locally presentable
categories. Imposing finite-limit equations amounts to passing to equifiers
of accessible functors. Locally presentable categories are closed under these
constructions, and the resulting forgetful functors are accessible.
\end{proof}

\begin{proposition}
\label{prop:mealy-category-locally-presentable}
The category
\[
\mathsf{Mealy}(\mathbb A,\mathbb B)
\]
is locally presentable.
\end{proposition}

\begin{proof}
The carrier category
\[
\mathcal E/Z
\]
is locally presentable because \(\mathcal E\) is locally presentable.

The assignment
\[
P\longmapsto A_1\times_{t_A,A_0,p}P
\]
is accessible, since it is built from pullback and dependent-sum functors
between locally presentable slices. The same is true for the finite iterated
pullback constructions appearing in the unit and multiplication laws.

Consider first the category of raw structures consisting of a carrier
\[
P\to Z,
\]
a transition map
\[
\delta:A_1\times_{t_A,A_0,p}P\to P,
\]
and an output map
\[
\omega:A_1\times_{t_A,A_0,p}P\to B_1
\]
satisfying the source-target typing equations. This is a category of
structures over \(\mathcal E/Z\) classified by accessible
dependent-polynomial operations.

The unit and multiplication laws for \(\delta\) and \(\omega\) are equations
between morphisms whose domains are obtained from \(P\) by finite pullback
constructions along fixed maps determined by \(\mathbb A\) and \(\mathbb B\).
Thus
\[
\mathsf{Mealy}(\mathbb A,\mathbb B)
\]
is obtained from a category of accessible dependent-polynomial structures by
imposing finite-limit equations. By
Lemma~\ref{lem:accessible-dependent-polynomial-structures}, it is locally
presentable.
\end{proof}

\begin{theorem}[Cofree Mealy systems]
\label{thm:cofree-mealy-systems}
The forgetful functor
\[
U_{\mathbb A,\mathbb B}:
\mathsf{Mealy}(\mathbb A,\mathbb B)
\to
\mathcal E/Z
\]
has a right adjoint
\[
R_{\mathbb A,\mathbb B}:
\mathcal E/Z
\to
\mathsf{Mealy}(\mathbb A,\mathbb B).
\]
\end{theorem}

\begin{proof}
By Proposition~\ref{prop:mealy-category-locally-presentable}, both
\[
\mathsf{Mealy}(\mathbb A,\mathbb B)
\qquad\text{and}\qquad
\mathcal E/Z
\]
are locally presentable. The functor
\[
U_{\mathbb A,\mathbb B}
\]
is accessible, since it is the forgetful functor from an accessible category
of structured objects over \(\mathcal E/Z\). By
Corollary~\ref{cor:mealy-forgetful-preserves-colimits}, it preserves
colimits. Therefore the adjoint functor theorem for locally presentable
categories gives a right adjoint
\[
R_{\mathbb A,\mathbb B}.
\]
\end{proof}

\begin{definition}[Cofree Mealy system]
\label{def:cofree-mealy-system}
For an object
\[
X\to Z
\]
of \(\mathcal E/Z\), the Mealy system
\[
R_{\mathbb A,\mathbb B}(X)
\]
is called the \textbf{cofree Mealy system} on \(X\).
\end{definition}

\begin{remark}
\label{rem:cofree-mealy-system-interpretation}
The counit of the adjunction
\[
U_{\mathbb A,\mathbb B}\dashv R_{\mathbb A,\mathbb B}
\]
at \(X\) is a carrier map
\[
U_{\mathbb A,\mathbb B}R_{\mathbb A,\mathbb B}(X)\to X.
\]
Thus \(R_{\mathbb A,\mathbb B}(X)\) is the universal Mealy system equipped
with an observation map into \(X\). The residual form of the cofree transpose
will be described after the final residual behaviour object has been
constructed explicitly.
\end{remark}

\section{The Mealy comonad}
\label{sec:mealy-comonad}

\begin{definition}[Mealy comonad]
\label{def:mealy-comonad}
The \textbf{Mealy comonad} associated to
\[
\mathbb A
\qquad\text{and}\qquad
\mathbb B
\]
is the comonad
\[
\mathbb M_{\mathbb A,\mathbb B}
:=
U_{\mathbb A,\mathbb B}R_{\mathbb A,\mathbb B}
:
\mathcal E/Z\to\mathcal E/Z
\]
induced by the adjunction
\[
U_{\mathbb A,\mathbb B}\dashv R_{\mathbb A,\mathbb B}.
\]
\end{definition}

\begin{theorem}[Mealy systems are coalgebras for the Mealy comonad]
\label{thm:mealy-systems-are-comonad-coalgebras}
There is an equivalence of categories
\[
\mathsf{Mealy}(\mathbb A,\mathbb B)
\simeq
\mathbf{Coalg}(\mathbb M_{\mathbb A,\mathbb B}).
\]
\end{theorem}

\begin{proof}
We use the comonadic Beck theorem. Since
\[
U_{\mathbb A,\mathbb B}
\]
has a right adjoint, it is comonadic if it reflects isomorphisms and creates
equalizers of \(U_{\mathbb A,\mathbb B}\)-split parallel pairs. We prove the
stronger statement that it creates all equalizers.

First, \(U_{\mathbb A,\mathbb B}\) reflects isomorphisms. Indeed, if a Mealy
transformation
\[
\theta:P\to Q
\]
is an isomorphism in \(\mathcal E/Z\), then its inverse automatically
preserves the Mealy transition and output maps, because the preservation
equations for \(\theta\) can be transported across the inverse.

Second, \(U_{\mathbb A,\mathbb B}\) creates equalizers. Let
\[
f,g:P\rightrightarrows Q
\]
be Mealy transformations, and let
\[
e:E\to P
\]
be their equalizer in \(\mathcal E/Z\). The transition map on \(P\) restricts
to \(E\). Indeed,
\[
f\delta_P(1_{A_1}\times e)
=
\delta_Q(1_{A_1}\times fe)
=
\delta_Q(1_{A_1}\times ge)
=
g\delta_P(1_{A_1}\times e).
\]
Hence \(\delta_P(1_{A_1}\times e)\) factors uniquely through \(E\). This
gives a transition map
\[
\delta_E:
A_1\times_{A_0}E\to E.
\]
The output map is
\[
\omega_E:=\omega_P(1_{A_1}\times e).
\]
The Mealy equations for \(E\) follow from those for \(P\). Thus the equalizer
is created by \(U_{\mathbb A,\mathbb B}\).

Therefore \(U_{\mathbb A,\mathbb B}\) is comonadic, and the comparison
functor
\[
\mathsf{Mealy}(\mathbb A,\mathbb B)
\to
\mathbf{Coalg}(\mathbb M_{\mathbb A,\mathbb B})
\]
is an equivalence.
\end{proof}

\section{Final behaviour}
\label{sec:final-mealy-behaviour}

Let
\[
1_Z:Z\to Z
\]
be the terminal object of the slice \(\mathcal E/Z\).

\begin{definition}[Behaviour system]
\label{def:behaviour-object}
The \textbf{final behaviour system} for Mealy morphisms
\[
\mathbb A\rightsquigarrow\mathbb B
\]
is
\[
\mathsf{Beh}_{\mathbb A,\mathbb B}
:=
R_{\mathbb A,\mathbb B}(1_Z).
\]
Its underlying carrier over
\[
Z=A_0\times B_0
\]
will also be denoted
\[
U_{\mathbb A,\mathbb B}\mathsf{Beh}_{\mathbb A,\mathbb B}
\]
when the distinction is needed.
\end{definition}

\begin{theorem}[Final behaviour coalgebra]
\label{thm:final-behaviour-coalgebra}
The Mealy system
\[
\mathsf{Beh}_{\mathbb A,\mathbb B}
\]
is terminal in
\[
\mathsf{Mealy}(\mathbb A,\mathbb B).
\]
Equivalently, the carrier
\[
U_{\mathbb A,\mathbb B}\mathsf{Beh}_{\mathbb A,\mathbb B}
=
\mathbb M_{\mathbb A,\mathbb B}(1_Z),
\]
equipped with its canonical Eilenberg--Moore coalgebra structure for the
comonad \(\mathbb M_{\mathbb A,\mathbb B}\), is terminal in
\[
\mathbf{Coalg}(\mathbb M_{\mathbb A,\mathbb B}).
\]
\end{theorem}

\begin{proof}
For every Mealy system \(P\), the cofree adjunction gives
\[
\mathsf{Mealy}(\mathbb A,\mathbb B)
(P,R_{\mathbb A,\mathbb B}(1_Z))
\cong
(\mathcal E/Z)(U_{\mathbb A,\mathbb B}P,1_Z).
\]
The object \(1_Z\) is terminal in \(\mathcal E/Z\), so the right-hand side is
a singleton. Hence
\[
R_{\mathbb A,\mathbb B}(1_Z)
\]
is terminal in \(\mathsf{Mealy}(\mathbb A,\mathbb B)\).
\end{proof}

\begin{definition}[Behaviour map]
\label{def:behaviour-map}
For a Mealy system \(P\), the unique Mealy transformation
\[
\beh_P:
P\to\mathsf{Beh}_{\mathbb A,\mathbb B}
\]
is called the \textbf{behaviour map} of \(P\).
\end{definition}

\subsection{Residual-slice form of behaviours}
\label{subsec:residual-slice-form-behaviours}

The final behaviour system can be read in a residual form. We now give an
explicit internal construction of this residual object.

Let
\[
v\in A_0
\]
be a generalized input object. The residual slice category
\[
\mathbb A/v
\]
has, in generalized-element notation, objects given by arrows
\[
r:u\to v
\]
of \(\mathbb A\). A morphism from
\[
r\circ a:w\to v
\]
to
\[
r:u\to v
\]
is given by an arrow
\[
a:w\to u
\]
of \(\mathbb A\). Thus a morphism is represented by a triangle
\[
w\xrightarrow{a}u\xrightarrow{r}v.
\]
The direction is from the longer residual history \(r\circ a\) to the shorter
residual history \(r\). This direction is forced by the right-action
orientation of Mealy morphisms.

\begin{definition}[Residual behaviour data]
\label{def:residual-behaviour-data}
A residual behaviour over
\[
(v,y)\in A_0\times B_0
\]
consists of the following data:
\begin{enumerate}
\item for every arrow \(r:u\to v\) of \(\mathbb A\), an object
\[
H(r)\in B_0;
\]
\item an equality
\[
H(1_v)=y;
\]
\item for every triangle
\[
w\xrightarrow{a}u\xrightarrow{r}v,
\]
an arrow
\[
H(a,r):H(r\circ a)\to H(r)
\]
in \(\mathbb B\);
\item identity and composition equations
\[
H(e_A(s_A r),r)=e_B(H(r)),
\]
and
\[
H(a\circ c,r)
=
H(a,r)\circ H(c,r\circ a)
\]
for every composable triple
\[
z\xrightarrow{c}w\xrightarrow{a}u\xrightarrow{r}v.
\]
\end{enumerate}
Equivalently, a residual behaviour over \((v,y)\) is a functor
\[
H:\mathbb A/v\to\mathbb B
\]
satisfying
\[
H(1_v)=y.
\]
\end{definition}

\begin{proposition}[Existence of the residual behaviour object]
\label{prop:residual-behaviour-object-exists}
Assume \(\mathcal E\) is locally cartesian closed. The residual behaviour data
of Definition~\ref{def:residual-behaviour-data} are represented by an object
\[
\operatorname{Res}_{\mathbb A,\mathbb B}\to A_0\times B_0.
\]
Thus generalized elements of the fibre over \((v,y)\) are precisely pointed
residual-slice functors
\[
H:\mathbb A/v\to\mathbb B,
\qquad
H(1_v)=y.
\]
\end{proposition}

\begin{proof}
The residual-arrow object over \(A_0\) is
\[
A_1\xrightarrow{t_A}A_0,
\]
whose fibre over \(v\) consists of arrows \(r:u\to v\). The residual-triangle
object over \(A_0\) is
\[
A_1\times_{s_A,A_0,t_A}A_1
\longrightarrow A_0,
\]
whose generalized elements are pairs
\[
(a,r)
\]
with
\[
w\xrightarrow{a}u\xrightarrow{r}v,
\]
projected to \(v=t_A(r)\). The residual-composable-triangle object is the
corresponding iterated pullback of triples
\[
(c,a,r)
\]
with
\[
z\xrightarrow{c}w\xrightarrow{a}u\xrightarrow{r}v.
\]

Because \(\mathcal E\) is locally cartesian closed, dependent products along
these indexing objects exist. They form the internal objects of assignments
\[
r\longmapsto H(r)
\]
and
\[
(a,r)\longmapsto H(a,r).
\]
The source-target conditions for \(H(a,r)\), the distinguished-point condition
\(H(1_v)=y\), and the identity and composition equations are equations between
maps obtained from these dependent products by finite-limit constructions.
Imposing them is therefore accomplished by finite limits, in particular by
equalizers.

The resulting object over \(A_0\times B_0\) represents exactly the data and
equations of Definition~\ref{def:residual-behaviour-data}.
\end{proof}

\begin{remark}[Why the residual functor is \(\mathbb B\)-valued]
\label{rem:residual-functor-variance}
Although Mealy executions are labelled by \(\mathbb B^{\mathrm{op}}\), the
residual-slice normal form is \(\mathbb B\)-valued. The reason is that
\(\mathbb A/v\) is oriented from longer residual histories to shorter ones:
a triangle
\[
w\xrightarrow{a}u\xrightarrow{r}v
\]
gives a morphism
\[
r\circ a\to r.
\]
Under a Mealy system, this residual morphism is sent to the output arrow
\[
q(\delta(r\circ a,x))\to q(\delta(r,x))
\]
in \(\mathbb B\). Thus the opposite variance has already been absorbed into
the definition of the residual slice.
\end{remark}

\begin{proposition}[Residual behaviours form the terminal Mealy system]
\label{prop:residual-behaviours-terminal}
The object
\[
\operatorname{Res}_{\mathbb A,\mathbb B}
\]
carries a canonical Mealy-system structure. It is terminal in
\[
\mathsf{Mealy}(\mathbb A,\mathbb B).
\]
Consequently there is a canonical isomorphism of Mealy systems
\[
\operatorname{Res}_{\mathbb A,\mathbb B}
\cong
\mathsf{Beh}_{\mathbb A,\mathbb B}.
\]
\end{proposition}

\begin{proof}
Let a residual behaviour be represented by
\[
H:\mathbb A/v\to\mathbb B,
\qquad
H(1_v)=y.
\]
For an admissible input arrow
\[
a:u\to v,
\]
define its derivative
\[
\partial_a H:\mathbb A/u\to\mathbb B
\]
by residual precomposition:
\[
(\partial_aH)(r)=H(a\circ r)
\]
for \(r:w\to u\). On morphisms, \(\partial_aH\) is defined by the same
precomposition operation on residual triangles. Its distinguished output
object is
\[
(\partial_aH)(1_u)=H(a).
\]
Thus the transition map sends
\[
(a,H)\longmapsto \partial_aH.
\]
The output arrow is
\[
H(a,1_v):H(a)\to H(1_v)=y.
\]
This defines the output component
\[
\omega_{\operatorname{Res}}(a,H):q(\partial_aH)\to q(H)
\]
in \(\mathbb B\).

The identity and multiplication laws follow from the identity and composition
laws for the functor \(H:\mathbb A/v\to\mathbb B\). Hence
\[
\operatorname{Res}_{\mathbb A,\mathbb B}
\]
is a Mealy system.

Now let \(P\) be any Mealy system. Define
\[
\beta_P:P\to\operatorname{Res}_{\mathbb A,\mathbb B}
\]
as follows. For a state \(x\in P\) with
\[
p(x)=v,
\qquad
q(x)=y,
\]
let
\[
H_x:\mathbb A/v\to\mathbb B
\]
be given on objects by
\[
H_x(r)=q(\delta_P(r,x))
\]
and on a residual triangle
\[
w\xrightarrow{a}u\xrightarrow{r}v
\]
by
\[
H_x(a,r)=\omega_P(a,\delta_P(r,x)).
\]
The Mealy typing equations give
\[
H_x(a,r):
H_x(r\circ a)\to H_x(r),
\]
and the Mealy unit and multiplication equations give functoriality. Moreover,
\[
H_x(1_v)=q(x)=y.
\]
Thus \(\beta_P\) is well-defined.

The equations
\[
\beta_P(\delta_P(a,x))=\partial_a(\beta_P(x))
\]
and
\[
\omega_P(a,x)=\omega_{\operatorname{Res}}(a,\beta_P(x))
\]
show that \(\beta_P\) is a Mealy transformation.

Finally, uniqueness follows because any Mealy transformation
\[
g:P\to\operatorname{Res}_{\mathbb A,\mathbb B}
\]
is forced, by transition preservation and preservation of the carrier map to
\(Z\), to satisfy
\[
g(x)(r)=q(\delta_P(r,x))
\]
on residual objects, and by output preservation to satisfy
\[
g(x)(a,r)=\omega_P(a,\delta_P(r,x))
\]
on residual arrows. Hence \(g=\beta_P\). Therefore
\[
\operatorname{Res}_{\mathbb A,\mathbb B}
\]
is terminal. Since
\[
\mathsf{Beh}_{\mathbb A,\mathbb B}
\]
is also terminal, the two are canonically isomorphic as Mealy systems.
\end{proof}

\begin{proposition}[Behaviours as pointed residual-slice functors]
\label{prop:behaviours-as-residual-slice-functors}
Let
\[
P
\]
be a Mealy system and let
\[
x\in P
\]
be a generalized state with
\[
p(x)=v,
\qquad
q(x)=y.
\]
The behaviour
\[
\beh_P(x)\in\mathsf{Beh}_{\mathbb A,\mathbb B}
\]
is represented by the functor
\[
H_x:\mathbb A/v\to\mathbb B
\]
whose value at the identity object
\[
1_v:v\to v
\]
is
\[
H_x(1_v)=y.
\]
It is defined on objects by
\[
H_x(r)=q(\delta_P(r,x))
\]
for each residual arrow
\[
r:u\to v.
\]
It is defined on morphisms by
\[
H_x(a:r\circ a\to r)
=
\omega_P(a,\delta_P(r,x)).
\]
The type is
\[
\omega_P(a,\delta_P(r,x)):
q(\delta_P(r\circ a,x))
\to
q(\delta_P(r,x)),
\]
so indeed
\[
H_x(a):H_x(r\circ a)\to H_x(r).
\]
The functoriality equations are exactly the Mealy unit and multiplication
equations:
\[
H_x(1_r)=1_{H_x(r)}
\]
and
\[
H_x(a\circ c)=H_x(a)\circ H_x(c).
\]
Conversely, the final behaviour system is the universal object of such
pointed residual-slice functors.
\end{proposition}

\begin{proof}
This is the concrete description of the terminal map
\[
P\to\operatorname{Res}_{\mathbb A,\mathbb B}
\]
from Proposition~\ref{prop:residual-behaviours-terminal}, transported across
the canonical isomorphism
\[
\operatorname{Res}_{\mathbb A,\mathbb B}
\cong
\mathsf{Beh}_{\mathbb A,\mathbb B}.
\]
\end{proof}

\begin{proposition}[Residual form of cofree transposes]
\label{prop:cofree-transpose-residual-form}
Let
\[
h:U_{\mathbb A,\mathbb B}P\to X
\]
be a map in \(\mathcal E/Z\), where \(P\) is a Mealy system. Let
\[
\widehat h:P\to R_{\mathbb A,\mathbb B}(X)
\]
be its transpose under
\[
U_{\mathbb A,\mathbb B}\dashv R_{\mathbb A,\mathbb B}.
\]
Then \(\widehat h\) is characterized by the following residual data, in the
sense that it is the unique Mealy transformation whose composite with the
cofree observation
\[
U_{\mathbb A,\mathbb B}R_{\mathbb A,\mathbb B}(X)\to X
\]
is \(h\).

For a generalized state \(x\in P\) with
\[
p(x)=v,
\]
and for every arrow
\[
r:u\to v
\]
of \(\mathbb A\), the residual state is
\[
\delta_P(r,x),
\]
and its residual observation is
\[
h(\delta_P(r,x)).
\]
For every further arrow
\[
a:w\to u,
\]
the residual one-step output is
\[
\omega_P(a,\delta_P(r,x)).
\]

These data satisfy the identity and composition equations inherited from the
Mealy laws. In particular,
\[
\delta_P(e_A(v),x)=x,
\]
and
\[
\delta_P(r\circ a,x)=\delta_P(a,\delta_P(r,x)).
\]
The output component satisfies
\[
\omega_P(e_A(u),\delta_P(r,x))
=
e_B(q(\delta_P(r,x))),
\]
and for
\[
r:u\to v,
\qquad
a:w\to u,
\qquad
c:z\to w,
\]
one has
\[
\omega_P(a\circ c,\delta_P(r,x))
=
\omega_P(a,\delta_P(r,x))
\circ
\omega_P(c,\delta_P(r\circ a,x)).
\]
\end{proposition}

\begin{proof}
The transpose \(\widehat h\) is the unique Mealy transformation into the
cofree system whose composite with the counit
\[
U_{\mathbb A,\mathbb B}R_{\mathbb A,\mathbb B}(X)\to X
\]
is \(h\). Since \(\widehat h\) is a Mealy transformation, it preserves
transitions and outputs. Repeated application of transition preservation says
that the image of \(x\) records the residual states
\[
\delta_P(r,x)
\]
for all admissible input arrows \(r\). Composing these residual states with
\(h\) gives the displayed residual observations, and output preservation gives
the displayed residual one-step outputs.

The equations are exactly the unit and multiplication laws for the transition
and output maps of the Mealy system \(P\). This proposition is a residual
description of the transpose of an already existing cofree map, not a
construction of \(R_{\mathbb A,\mathbb B}(X)\).
\end{proof}

\begin{remark}
\label{rem:residual-slice-behaviour-meaning}
This proposition is the general internal version of the statement that a
Mealy behaviour records all future outputs. Here future inputs are not
necessarily words over an alphabet. They are arrows in the input internal
category, and the residual futures at an input object \(v\) are organized by
the slice category
\[
\mathbb A/v.
\]
Thus the correct general behaviour object is an object of pointed
residual-slice functors, not an arbitrary exponential of words into outputs.
\end{remark}

\section{Derivatives}
\label{sec:mealy-behaviour-derivatives}

Since
\[
\mathsf{Beh}_{\mathbb A,\mathbb B}
\]
is itself a Mealy system, it has a transition map
\[
\delta_{\mathrm{Beh}}:
A_1\times_{t_A,A_0,p_{\mathrm{Beh}}}
\mathsf{Beh}_{\mathbb A,\mathbb B}
\to
\mathsf{Beh}_{\mathbb A,\mathbb B}.
\]

\begin{definition}[Derivative]
\label{def:behaviour-derivative}
For an admissible input arrow
\[
a:u\to v
\]
and a behaviour \(\beta\) lying over \(v\), define
\[
\partial_a\beta
:=
\delta_{\mathrm{Beh}}(a,\beta).
\]
\end{definition}

\begin{proposition}[Derivative laws]
\label{prop:derivative-laws}
The derivatives in
\[
\mathsf{Beh}_{\mathbb A,\mathbb B}
\]
satisfy
\[
\partial_{e_A(v)}\beta=\beta
\]
and, for composable arrows
\[
a:u\to v,
\qquad
b:v\to w,
\]
and \(\beta\) over \(w\),
\[
\partial_{b\circ a}\beta
=
\partial_a(\partial_b\beta).
\]
\end{proposition}

\begin{proof}
These are exactly the unit and multiplication laws for the Mealy transition
map of the final behaviour system.
\end{proof}

\begin{proposition}[Derivatives as residual precomposition]
\label{prop:derivatives-as-residual-precomposition}
Let a behaviour
\[
\beta\in\mathsf{Beh}_{\mathbb A,\mathbb B}
\]
over \(v\) be represented by a pointed residual-slice functor
\[
H:\mathbb A/v\to\mathbb B.
\]
Let
\[
r:u\to v
\]
be an admissible input arrow. There is a postcomposition functor
\[
r_\ast:\mathbb A/u\to\mathbb A/v
\]
defined on objects by
\[
a:w\to u
\quad\longmapsto\quad
r\circ a:w\to v.
\]
Then the derivative of \(H\) by \(r\) is
\[
\partial_r H
=
H\circ r_\ast:
\mathbb A/u\to\mathbb B.
\]
In particular,
\[
(\partial_rH)(1_u)=H(r).
\]
Thus the output interface of the derivative is the residual output interface
of the original behaviour at \(r\).
\end{proposition}

\begin{proof}
A behaviour is a state of the final Mealy system itself. Apply
Proposition~\ref{prop:behaviours-as-residual-slice-functors} to
\[
P=\mathsf{Beh}_{\mathbb A,\mathbb B}
\]
and \(x=\beta\). The derivative by \(r:u\to v\) is the residual state
\[
\delta_{\mathrm{Beh}}(r,\beta).
\]
For any arrow
\[
a:w\to u,
\]
the Mealy multiplication law gives
\[
\delta_{\mathrm{Beh}}(a,\delta_{\mathrm{Beh}}(r,\beta))
=
\delta_{\mathrm{Beh}}(r\circ a,\beta).
\]
Therefore the residual-slice functor of
\[
\partial_r\beta
\]
takes \(a\) to the value of \(H\) at \(r\circ a\). The same computation on
residual triangles gives equality on arrows. Hence
\[
\partial_r H=H\circ r_\ast.
\]
\end{proof}

\begin{remark}
\label{rem:derivatives-as-residual-restriction}
Derivation is not an additional operation imposed on behaviours. It is
restriction of a residual-slice functor along the functor that extends a
future input history by a fixed prefix. This is the abstract form of the usual
residual operation on languages or stream behaviours.
\end{remark}

\begin{proposition}[Behaviour maps preserve derivatives and outputs]
\label{prop:behaviour-maps-preserve-derivatives}
For every Mealy system \(P\), every admissible input arrow \(a\), and every
state \(x\), one has
\[
\beh_P(\delta_P(a,x))
=
\partial_a(\beh_P(x)).
\]
Moreover,
\[
\omega_P(a,x)
=
\omega_{\mathrm{Beh}}(a,\beh_P(x)).
\]
\end{proposition}

\begin{proof}
The behaviour map
\[
\beh_P:P\to\mathsf{Beh}_{\mathbb A,\mathbb B}
\]
is a Mealy transformation. The displayed equations are exactly preservation
of the transition and output components.
\end{proof}

\section{Behavioural equivalence and bisimulation}
\label{sec:behavioural-equivalence-bisimulation}

\begin{definition}[Behavioural equivalence]
\label{def:behavioural-equivalence}
For a Mealy system \(P\), define its behavioural equivalence relation by the
kernel pair
\[
{\sim_{\beh}}
:=
P\times_{\mathsf{Beh}_{\mathbb A,\mathbb B}}P
\]
of the behaviour map
\[
\beh_P:P\to\mathsf{Beh}_{\mathbb A,\mathbb B}.
\]
\end{definition}

\begin{definition}[Mealy bisimulation]
\label{def:mealy-bisimulation}
Let \(P\) be a Mealy system with carrier
\[
P\to Z=A_0\times B_0.
\]
A \textbf{Mealy bisimulation} on \(P\) is a subobject
\[
R\hookrightarrow P\times_ZP
\]
which carries a Mealy-system structure for which the two projections
\[
R\rightrightarrows P
\]
are Mealy transformations.

Equivalently, in generalized-element notation, for every
\[
(x,x')\in R
\]
and every admissible input arrow
\[
a:u\to p(x)=p(x'),
\]
one has
\[
(\delta_P(a,x),\delta_P(a,x'))\in R,
\]
and, under the induced equality of output interfaces, one has
\[
\omega_P(a,x)=\omega_P(a,x').
\]
The output map on \(R\) is then the common value of these two outputs.
\end{definition}

\begin{theorem}[Behavioural equivalence is largest bisimulation]
\label{thm:behavioural-equivalence-largest-bisimulation}
The relation
\[
{\sim_{\beh}}
\]
is the largest Mealy bisimulation on \(P\).
\end{theorem}

\begin{proof}
The kernel pair
\[
P\times_{\mathsf{Beh}_{\mathbb A,\mathbb B}}P
\]
inherits a Mealy-system structure by the universal property of the pullback,
because the two projections to \(P\) have equal composites with the behaviour
map and the behaviour map is a Mealy transformation. Explicitly, if two
states \(x,x'\) have the same behaviour, then, for every admissible input
arrow \(a\),
\[
\beh_P(\delta_P(a,x))
=
\partial_a(\beh_P(x))
=
\partial_a(\beh_P(x'))
=
\beh_P(\delta_P(a,x')).
\]
Thus the derivatives are again behaviourally equivalent. Moreover,
\[
\omega_P(a,x)
=
\omega_{\mathrm{Beh}}(a,\beh_P(x))
=
\omega_{\mathrm{Beh}}(a,\beh_P(x'))
=
\omega_P(a,x').
\]
Hence \({\sim_{\beh}}\) is a bisimulation.

Conversely, let
\[
R\hookrightarrow P\times_ZP
\]
be any Mealy bisimulation. Equivalently, the projections
\[
\pi_1,\pi_2:R\rightrightarrows P
\]
are Mealy transformations. Therefore the two composites
\[
R\rightrightarrows P\xrightarrow{\beh_P}
\mathsf{Beh}_{\mathbb A,\mathbb B}
\]
are two Mealy transformations into the terminal Mealy system. They are
therefore equal. Hence \(R\) factors through the kernel pair of
\[
\beh_P:P\to\mathsf{Beh}_{\mathbb A,\mathbb B}.
\]
Thus \(R\) is contained in \({\sim_{\beh}}\), so \({\sim_{\beh}}\) is the
largest bisimulation.
\end{proof}

\section{The one-object monoid case}
\label{sec:one-object-monoid-case}

This section records the explicit form of the preceding construction when
\[
\mathcal E=\mathbf{Set}
\]
and
\[
\mathbb A=M,
\qquad
\mathbb B=N
\]
are one-object internal categories, hence monoids.

We keep the convention
\[
mn=m\circ n.
\]
Thus \(mn\) means first \(n\), then \(m\), as categorical composition. Since
the Mealy input action is a right action, the operational order is reversed:
the state action by \(mn\) processes \(m\) first and \(n\) second.

A Mealy system is a right \(M\)-set \(P\) equipped with an \(N\)-valued cocycle
\[
\omega:M\times P\to N.
\]
The laws are
\[
x\cdot 1=x,
\qquad
x\cdot(mn)=(x\cdot m)\cdot n,
\]
and
\[
\omega(1,x)=1,
\qquad
\omega(mn,x)=\omega(m,x)\omega(n,x\cdot m).
\]

\begin{proposition}[Cofree monoid-input Mealy systems]
\label{prop:cofree-monoid-mealy}
For a set \(X\), the underlying carrier of the cofree \((M,N)\)-Mealy system
on \(X\) is
\[
\mathbb M_{M,N}(X)
=
\left\{
(\ell,\lambda)
\ \middle|\
\ell:M\to X,\
\lambda:M\times M\to N,\
\lambda(r,1)=1,\
\lambda(r,mn)=\lambda(r,m)\lambda(rm,n)
\right\}.
\]
The observation map is
\[
\varepsilon_X(\ell,\lambda)=\ell(1).
\]
The right \(M\)-action is
\[
(\ell,\lambda)\cdot m=(\ell_m,\lambda_m),
\]
where
\[
\ell_m(r)=\ell(mr),
\qquad
\lambda_m(r,n)=\lambda(mr,n).
\]
The output is
\[
\omega_{\mathbb M}(m,(\ell,\lambda))=\lambda(1,m).
\]
\end{proposition}

\begin{proof}
The displayed equations for \(\lambda\) imply
\[
\lambda_m(r,1)=\lambda(mr,1)=1
\]
and
\[
\lambda_m(r,ab)
=
\lambda(mr,ab)
=
\lambda(mr,a)\lambda(mra,b)
=
\lambda_m(r,a)\lambda_m(ra,b).
\]
Thus residuals are again valid states.

The action laws follow from associativity in \(M\):
\[
((\ell,\lambda)\cdot m)\cdot n
=
(\ell,\lambda)\cdot(mn),
\qquad
(\ell,\lambda)\cdot1=(\ell,\lambda).
\]
The output unit law is
\[
\omega_{\mathbb M}(1,(\ell,\lambda))=\lambda(1,1)=1.
\]
The cocycle law is
\[
\omega_{\mathbb M}(mn,(\ell,\lambda))
=
\lambda(1,mn)
=
\lambda(1,m)\lambda(m,n)
=
\omega_{\mathbb M}(m,(\ell,\lambda))
\omega_{\mathbb M}(n,(\ell,\lambda)\cdot m).
\]
Hence the displayed carrier is a Mealy system.

Given any \((M,N)\)-Mealy system \(P\) and any function
\[
h:P\to X,
\]
define
\[
\widehat h:P\to\mathbb M_{M,N}(X)
\]
by
\[
\widehat h(x)=(\ell_x,\lambda_x),
\]
where
\[
\ell_x(r)=h(x\cdot r),
\qquad
\lambda_x(r,m)=\omega(m,x\cdot r).
\]
The action and cocycle laws for \(P\) imply that \(\lambda_x\) satisfies the
required equations. The map \(\widehat h\) is a Mealy transformation, and
\[
\varepsilon_X\widehat h=h.
\]

For uniqueness, let
\[
g:P\to\mathbb M_{M,N}(X)
\]
be a Mealy transformation satisfying
\[
\varepsilon_Xg=h.
\]
If
\[
g(x)=(\ell,\lambda),
\]
then preservation of the action gives
\[
\ell(r)
=
\varepsilon_X(g(x)\cdot r)
=
h(x\cdot r),
\]
and preservation of outputs gives
\[
\lambda(r,m)
=
\omega_{\mathbb M}(m,g(x)\cdot r)
=
\omega(m,x\cdot r).
\]
Thus \(g=\widehat h\). Therefore the displayed system is cofree.
\end{proof}

\begin{remark}
\label{rem:monoid-cofree-interpretation}
In this formula, \(\ell(r)\) is the residual observation after processing the
input monoid element \(r\), while \(\lambda(r,m)\) is the output produced by
processing \(m\) after residual history \(r\). The equation
\[
\lambda(r,mn)=\lambda(r,m)\lambda(rm,n)
\]
says that output along a composite input \(mn\) is obtained by first outputting
along \(m\), then outputting along \(n\) at the residual position \(rm\).
\end{remark}

\begin{corollary}[Final monoid-input behaviour]
\label{cor:final-monoid-behaviour}
The final \((M,N)\)-Mealy behaviour object is
\[
\mathsf{Beh}_{M,N}
=
\left\{
\lambda:M\times M\to N
\ \middle|\
\lambda(r,1)=1,\
\lambda(r,mn)=\lambda(r,m)\lambda(rm,n)
\right\}.
\]
For a Mealy system \(P\), the behaviour map sends
\[
x\in P
\]
to
\[
\lambda_x(r,m)=\omega(m,x\cdot r).
\]
\end{corollary}

\begin{proposition}[One-object residual-slice form]
\label{prop:one-object-residual-slice-form}
Let
\[
M/\ast
\]
be the slice category of the one-object category associated to \(M\). Its
objects are elements
\[
r\in M,
\]
and a morphism
\[
rm\to r
\]
is determined by an element
\[
m\in M.
\]
Then
\[
\mathsf{Beh}_{M,N}
\cong
[M/\ast,N],
\]
where \(N\) is regarded as a one-object category. Under this isomorphism, a
functor
\[
M/\ast\to N
\]
corresponds to the cocycle
\[
\lambda:M\times M\to N
\]
given by the value of the functor on the morphism
\[
rm\to r.
\]
Here multiplication in \(N\) is written with the same convention as
categorical composition in the one-object category \(N\).
\end{proposition}

\begin{proof}
A functor \(M/\ast\to N\) assigns to each morphism
\[
rm\to r
\]
an element
\[
\lambda(r,m)\in N.
\]
Identity preservation gives
\[
\lambda(r,1)=1.
\]
Composition preservation gives
\[
\lambda(r,mn)=\lambda(r,m)\lambda(rm,n).
\]
These are precisely the equations defining \(\mathsf{Beh}_{M,N}\). Conversely,
any such cocycle defines a functor by the displayed assignment.
\end{proof}

\begin{example}[Derivatives in the monoid case]
\label{ex:monoid-derivatives}
For a behaviour
\[
\lambda\in\mathsf{Beh}_{M,N}
\]
and \(m\in M\), the derivative is
\[
(\partial_m\lambda)(r,n)=\lambda(mr,n).
\]
More generally,
\[
(\partial_r\lambda)(s,m)=\lambda(rs,m).
\]
For a system \(P\),
\[
\beh_P(x\cdot m)=\partial_m(\beh_P(x)).
\]
\end{example}

\begin{corollary}[Free-input recovery of the classical behaviour object]
\label{cor:free-input-classical-behaviour-object}
Suppose the input monoid is the free monoid
\[
M=I^\ast
\]
on a set \(I\). Then restriction of a coherent residual cocycle to generator
inputs gives a canonical isomorphism
\[
\mathsf{Beh}_{I^\ast,N}
\cong
N^{I^+},
\]
where \(I^+\) is the set of nonempty words. Equivalently,
\[
I^+\cong I^\ast\times I,
\qquad
(r,i)\longmapsto ri,
\]
by taking the last letter of a nonempty word. Here \(ri\) denotes ordinary
word notation in operational order: the residual history \(r\) followed by
the generator \(i\).

Given a cocycle
\[
\lambda:I^\ast\times I^\ast\to N,
\]
define
\[
\beta_\lambda:I^+\to N
\]
by
\[
\beta_\lambda(ri)=\lambda(r,i),
\qquad
r\in I^\ast,\ i\in I.
\]
Conversely, given
\[
\beta:I^+\to N,
\]
define
\[
\lambda_\beta(r,\epsilon)=1,
\]
and for
\[
w=i_1\cdots i_n,
\qquad n\ge1,
\]
define
\[
\lambda_\beta(r,w)
=
\beta(ri_1)
\beta(ri_1i_2)
\cdots
\beta(ri_1\cdots i_n).
\]
These assignments are inverse.
\end{corollary}

\begin{proof}
The cocycle equation implies by induction on word length that
\[
\lambda(r,i_1\cdots i_n)
=
\lambda(r,i_1)
\lambda(ri_1,i_2)
\cdots
\lambda(ri_1\cdots i_{n-1},i_n).
\]
Hence \(\lambda\) is uniquely determined by its values on pairs
\[
(r,i)\in I^\ast\times I.
\]
Conversely, the displayed product formula defines a cocycle because
concatenation in \(I^\ast\) is free and multiplication in \(N\) is
associative. The unit condition is immediate from
\[
\lambda_\beta(r,\epsilon)=1.
\]
Thus the two constructions are inverse.
\end{proof}

\begin{corollary}[Derivatives in the classical normal form]
\label{cor:free-input-derivatives-classical-normal-form}
Under the isomorphism
\[
\mathsf{Beh}_{I^\ast,N}\cong N^{I^+},
\]
the derivative by a word
\[
r\in I^\ast
\]
is the residual shift
\[
(\partial_r\beta)(w)=\beta(rw),
\qquad
w\in I^+.
\]
In particular, for a generator
\[
i\in I,
\]
one has
\[
(\partial_i\beta)(w)=\beta(iw).
\]
Although the generator-level normal form writes a nonempty word by its last
letter,
\[
ri\in I^+,
\]
the derivative operation is the left residual shift
\[
w\longmapsto rw.
\]
This is a consequence of the right-action convention and the multiplication
convention \(mn=m\circ n\).
\end{corollary}

\begin{remark}[Agreement with the standard final Mealy coalgebra]
\label{rem:agreement-standard-final-mealy-coalgebra}
The native one-object behaviour object for this chapter is the coherent
cocycle object
\[
\mathsf{Beh}_{M,N}
\hookrightarrow
N^{M\times M}.
\]
When \(M=I^\ast\) is free and one restricts the \(I^\ast\)-input Mealy
structure to generator inputs \(i\in I\), the final behaviour object becomes
\[
N^{I^+}.
\]
This is the standard final coalgebra shape for the classical deterministic
Mealy functor
\[
X\mapsto (N\times X)^I.
\]
Equivalently, it is the function-set form of the final coalgebra often
described by causal stream functions
\[
I^\omega\to N^\omega
\]
\cite{RuttenUniversalCoalgebra,SilvaCoalgebraicAutomata,PattinsonCoalgebraNotes}.
The point of Corollary~\ref{cor:free-input-classical-behaviour-object} is that
the classical object appears as the free-input normal form of the coherent
monoid-input behaviour object.
\end{remark}

\begin{example}[Behavioural equivalence in the monoid case]
\label{ex:monoid-behavioural-equivalence}
For an \((M,N)\)-Mealy system \(P\),
\[
x\sim_{\beh}x'
\]
if and only if
\[
\omega(m,x\cdot r)=\omega(m,x'\cdot r)
\]
for all
\[
r,m\in M.
\]
Thus behavioural equivalence is equality of residual output behaviour.
\end{example}

\section{Regular realization theory}
\label{sec:regular-realization-theory}

For this section, assume in addition that
\[
\mathcal E
\]
is regular. Then every slice
\[
\mathcal E/Z
\]
is regular, and images below are taken in the slice unless explicitly stated
otherwise.

\begin{definition}[Realized behaviour object]
\label{def:realized-behaviour-object}
For a Mealy system \(P\), define
\[
\operatorname{Real}(P)
\]
to be the image in \(\mathcal E/Z\) of the behaviour map
\[
\beh_P:P\to\mathsf{Beh}_{\mathbb A,\mathbb B}.
\]
Thus there is a regular epi--mono factorization in \(\mathcal E/Z\)
\[
P\twoheadrightarrow
\operatorname{Real}(P)
\hookrightarrow
\mathsf{Beh}_{\mathbb A,\mathbb B}.
\]
\end{definition}

\begin{proposition}
\label{prop:realized-behaviour-submealy}
The object
\[
\operatorname{Real}(P)
\]
is a sub-Mealy system of
\[
\mathsf{Beh}_{\mathbb A,\mathbb B}.
\]
\end{proposition}

\begin{proof}
Let
\[
P\twoheadrightarrow I\hookrightarrow\mathsf{Beh}_{\mathbb A,\mathbb B}
\]
be the image factorization of \(\beh_P\) in \(\mathcal E/Z\). Since regular
epis are stable under pullback, the induced map
\[
A_1\times_{A_0}P
\twoheadrightarrow
A_1\times_{A_0}I
\]
is a regular epi.

The derivative map of the final behaviour system sends
\[
A_1\times_{A_0}I
\]
back into \(I\). To see this, pull back along the regular epi above. The
resulting composite agrees with the image of
\[
\delta_P:
A_1\times_{A_0}P\to P
\]
followed by the image map
\[
P\twoheadrightarrow I.
\]
In a regular category, pullback along a regular epimorphism reflects inclusion
of subobjects. Since after pulling back along the regular epimorphism the
derivative lands in \(I\), it already lands in \(I\) before pullback. Hence
the derivative map restricts to \(I\).

The output map restricts by restriction of the final-system output map to
\[
A_1\times_{A_0}I.
\]
The Mealy equations hold because they hold in the final system. Hence \(I\) is
a sub-Mealy system.
\end{proof}

\begin{example}[Realized behaviours in the monoid case]
\label{ex:monoid-realized-behaviours}
For an \((M,N)\)-Mealy system \(P\),
\[
\operatorname{Real}(P)
=
\{\lambda_x\mid x\in P\}
\subseteq
\mathsf{Beh}_{M,N}.
\]
It is closed under derivatives because
\[
\partial_m\lambda_x=\lambda_{x\cdot m}.
\]
\end{example}

\section{Derivative closure of a specification}
\label{sec:derivative-closure-specification}

Assume throughout this section that
\[
\mathcal E
\]
is regular.

Let
\[
k:K\to\mathsf{Beh}_{\mathbb A,\mathbb B}
\]
be a behaviour specification in \(\mathcal E/Z\). Write
\[
p_K:=p_{\mathrm{Beh}}k,
\qquad
q_K:=q_{\mathrm{Beh}}k.
\]

Define the residual-domain object
\[
A_1\ltimes K
:=
A_1\times_{t_A,A_0,p_K}K.
\]
The derivative evaluation map is
\[
d_K:
A_1\ltimes K
\to
\mathsf{Beh}_{\mathbb A,\mathbb B}
\]
given by
\[
d_K(a,k)=\partial_a k.
\]
We regard
\[
A_1\ltimes K
\]
as an object of \(\mathcal E/Z\) through the composite
\[
A_1\ltimes K
\xrightarrow{d_K}
\mathsf{Beh}_{\mathbb A,\mathbb B}
\to
Z.
\]
With this slice structure, \(d_K\) is a morphism in \(\mathcal E/Z\). In
generalized-element notation, the induced base point is
\[
(a,k)\longmapsto
\bigl(s_A(a),q_{\mathrm{Beh}}(\partial_a k)\bigr).
\]

\begin{definition}[Derivative closure]
\label{def:derivative-closure}
Define
\[
\operatorname{Der}(K)
\]
to be the image in \(\mathcal E/Z\) of
\[
d_K:
A_1\ltimes K
\to
\mathsf{Beh}_{\mathbb A,\mathbb B}.
\]
Thus
\[
A_1\ltimes K
\twoheadrightarrow
\operatorname{Der}(K)
\hookrightarrow
\mathsf{Beh}_{\mathbb A,\mathbb B}
\]
is the regular epi--mono factorization of \(d_K\) in \(\mathcal E/Z\).
\end{definition}

\begin{remark}
\label{rem:derivative-closure-full-arrow-object}
Here \(A_1\) is the full arrow object of the internal category \(\mathbb A\),
not a chosen generator object. Thus \(d_K\) evaluates all derivatives by all
arrows of \(\mathbb A\). If one instead worked with a generator subobject of
\(A_1\), the derivative closure would have to be defined by an iterative image
construction.
\end{remark}

\begin{remark}[Derivative closure as residual realization]
\label{rem:derivative-closure-as-residual-realization}
The object
\[
\operatorname{Der}(K)
\]
is the residual-state system generated by the specified behaviours \(K\).
Since \(A_1\) is the full arrow object of the input internal category, the map
\[
A_1\ltimes K\to\mathsf{Beh}_{\mathbb A,\mathbb B}
\]
already evaluates all residual derivatives at once. Thus
\(\operatorname{Der}(K)\) is not obtained by repeatedly closing under a chosen
generator set. It is the image of the full residual-evaluation map.
\end{remark}

\begin{theorem}[Derivative closure]
\label{thm:derivative-closure}
The object
\[
\operatorname{Der}(K)
\]
is the smallest sub-Mealy system of
\[
\mathsf{Beh}_{\mathbb A,\mathbb B}
\]
containing the image of \(K\).
\end{theorem}

\begin{proof}
First, \(\operatorname{Der}(K)\) contains the image of
\[
k:K\to\mathsf{Beh}_{\mathbb A,\mathbb B}.
\]
Indeed, the unit map
\[
K\to A_1\ltimes K,
\qquad
k\longmapsto (e_A(p_Kk),k),
\]
satisfies
\[
d_K(e_A(p_Kk),k)=k.
\]
Hence \(k\) factors through the image of \(d_K\).

Second, \(\operatorname{Der}(K)\) is closed under derivatives. Let
\[
A_1\ltimes K
\twoheadrightarrow
\operatorname{Der}(K)
\hookrightarrow
\mathsf{Beh}_{\mathbb A,\mathbb B}
\]
be the image factorization of \(d_K\). Since regular epimorphisms are stable
under pullback, the induced map
\[
A_1\times_{t_A,A_0,p_D}(A_1\ltimes K)
\twoheadrightarrow
A_1\times_{t_A,A_0,p_{\operatorname{Der}}}
\operatorname{Der}(K)
\]
is a regular epi, where
\[
p_D
=
p_{\mathrm{Beh}}d_K
=
s_A\pi_1
\]
is the \(A_0\)-component of the derivative \(d_K(a,k)=\partial_a k\).

A generalized element of the left-hand side is a triple
\[
(c,a,k)
\]
with
\[
a:u\to v,
\qquad
c:r\to u,
\qquad
p_K(k)=v.
\]
After pulling back along the regular epimorphism above, the composite
derivative map is given by
\[
(c,a,k)\longmapsto \partial_c(\partial_a k).
\]
By Proposition~\ref{prop:derivative-laws},
\[
\partial_c(\partial_a k)
=
\partial_{a\circ c}k.
\]
The internal composition map gives
\[
m_A(c,a)=a\circ c:r\to v.
\]
Hence the pulled-back composite factors through \(d_K\) via the map
\[
(c,a,k)\longmapsto (m_A(c,a),k).
\]
It therefore factors through
\[
\operatorname{Der}(K)\hookrightarrow\mathsf{Beh}_{\mathbb A,\mathbb B}.
\]
Pullback along a regular epimorphism reflects inclusion of subobjects in a
regular category. Hence the original derivative map
\[
A_1\times_{A_0}\operatorname{Der}(K)
\to
\mathsf{Beh}_{\mathbb A,\mathbb B}
\]
factors through \(\operatorname{Der}(K)\). Therefore
\(\operatorname{Der}(K)\) is closed under derivatives.

Finally, any sub-Mealy system of \(\mathsf{Beh}_{\mathbb A,\mathbb B}\)
containing the image of \(K\) must contain every derivative of every element
of \(K\). Hence it must contain the image of \(d_K\). This proves the
minimality claim.
\end{proof}

\begin{remark}[Minimality before quotienting]
\label{rem:minimality-before-quotienting}
The derivative closure is a sub-Mealy system of the final behaviour object.
Hence it is already behaviourally separated: two different points of
\[
\operatorname{Der}(K)
\]
are different behaviours.

Thus
\[
\operatorname{Der}(K)
\]
is the minimal separated realization generated by the specification. The
Nerode quotient theorem explains how to obtain this separated realization
from raw residual syntax
\[
A_1\ltimes K
\]
by identifying derivatives with the same behaviour.
\end{remark}

\begin{example}[Derivative closure in the monoid case]
\label{ex:monoid-derivative-closure}
For a behaviour
\[
\lambda\in\mathsf{Beh}_{M,N},
\]
the derivative closure is the image of
\[
M\to\mathsf{Beh}_{M,N},
\qquad
r\longmapsto\partial_r\lambda.
\]
Thus
\[
\operatorname{Der}(\lambda)
=
\{\partial_r\lambda\mid r\in M\}.
\]
This is the residual-state system generated by \(\lambda\).
\end{example}

\section{The Nerode quotient theorem}
\label{sec:nerode-quotient-theorem}

For this section, assume in addition that
\[
\mathcal E
\]
is Barr-exact. Then every slice
\[
\mathcal E/Z
\]
is Barr-exact.

Let
\[
k:K\to\mathsf{Beh}_{\mathbb A,\mathbb B}
\]
be a behaviour specification.

\begin{definition}[Nerode equivalence]
\label{def:nerode-equivalence}
The \textbf{Nerode equivalence} associated to \(k\) is the kernel pair in
\(\mathcal E/Z\) of
\[
d_K:
A_1\ltimes K
\to
\mathsf{Beh}_{\mathbb A,\mathbb B}.
\]
Thus two generalized derivatives
\[
(a,k)
\qquad\text{and}\qquad
(a',k')
\]
are Nerode equivalent precisely when
\[
\partial_a k=\partial_{a'}k'.
\]
\end{definition}

\begin{remark}[What the Nerode equivalence identifies]
\label{rem:what-nerode-identifies}
The object
\[
A_1\ltimes K
\]
is the object of generalized residual expressions. A point
\[
(a,k)
\]
means: start from the specified behaviour \(k\), then take the derivative by
the input arrow \(a\).

The Nerode equivalence identifies two such generalized residuals precisely
when they determine the same behaviour:
\[
(a,k)\equiv_K(a',k')
\quad\Longleftrightarrow\quad
\partial_a k=\partial_{a'}k'.
\]
Thus the Nerode quotient is not a quotient of an arbitrary state space. It is
the quotient of the freely generated residual descriptions by behavioural
equality.
\end{remark}

\begin{theorem}[General Nerode realization theorem]
\label{thm:general-nerode-realization}
There is a canonical isomorphism in \(\mathcal E/Z\)
\[
(A_1\ltimes K)/{\equiv_K}
\cong
\operatorname{Der}(K),
\]
where \(\equiv_K\) denotes the Nerode equivalence. Transporting the
sub-Mealy-system structure of \(\operatorname{Der}(K)\) across this
isomorphism equips the quotient
\[
(A_1\ltimes K)/{\equiv_K}
\]
with a canonical Mealy-system structure. With this structure, the isomorphism
is an isomorphism of Mealy systems. Consequently,
\[
\text{Nerode quotient}
=
\text{derivative closure}.
\]
\end{theorem}

\begin{proof}
By definition, \(\operatorname{Der}(K)\) is the image of
\[
d_K:
A_1\ltimes K
\to
\mathsf{Beh}_{\mathbb A,\mathbb B}
\]
in \(\mathcal E/Z\). Factor \(d_K\) as
\[
A_1\ltimes K
\twoheadrightarrow
\operatorname{Der}(K)
\hookrightarrow
\mathsf{Beh}_{\mathbb A,\mathbb B}.
\]
The kernel pair of \(d_K\) is the kernel pair of the regular epimorphism
\[
A_1\ltimes K
\twoheadrightarrow
\operatorname{Der}(K).
\]
Since \(\mathcal E/Z\) is Barr-exact, this regular epimorphism is the
effective quotient of its kernel pair. Therefore
\[
(A_1\ltimes K)/{\equiv_K}
\cong
\operatorname{Der}(K)
\]
in \(\mathcal E/Z\).

By Theorem~\ref{thm:derivative-closure}, \(\operatorname{Der}(K)\) is a
sub-Mealy system of the final behaviour system. Transporting this Mealy
structure across the displayed isomorphism gives the quotient its canonical
Mealy-system structure, and the isomorphism is then an isomorphism of Mealy
systems by construction.
\end{proof}

\begin{remark}[Quotient form of the residual realization]
\label{rem:quotient-form-residual-realization}
The theorem says that the two standard constructions of the minimal residual
system agree:
\[
\operatorname{Der}(K)
=
\operatorname{im}
\left(
A_1\ltimes K\to\mathsf{Beh}_{\mathbb A,\mathbb B}
\right)
\]
and
\[
\operatorname{Der}(K)
\cong
(A_1\ltimes K)/{\equiv_K}.
\]
The image description views the minimal realization as a subobject of the
final behaviour object. The quotient description views it as residual syntax
modulo behavioural equality. Barr-exactness is precisely what guarantees that
these two descriptions coincide.
\end{remark}

\begin{corollary}
\label{cor:nerode-minimal-realization}
The derivative closure
\[
\operatorname{Der}(K)
\]
is the canonical minimal realization of the specification \(K\).
\end{corollary}

\begin{proof}
By Theorem~\ref{thm:derivative-closure}, \(\operatorname{Der}(K)\) is the
smallest sub-Mealy system of the final behaviour object containing the image
of \(K\). By Theorem~\ref{thm:general-nerode-realization}, it is also the
effective quotient of the generalized derivative object by the Nerode
equivalence. Thus it is simultaneously the derivative closure and the Nerode
quotient.
\end{proof}

\begin{example}[Nerode quotient in the monoid case]
\label{ex:monoid-nerode-quotient}
For a behaviour
\[
\lambda\in\mathsf{Beh}_{M,N},
\]
define
\[
r\equiv_\lambda s
\]
if and only if
\[
\partial_r\lambda=\partial_s\lambda.
\]
Equivalently,
\[
r\equiv_\lambda s
\]
if and only if, for all \(u,m\in M\),
\[
\lambda(ru,m)=\lambda(su,m).
\]
Then
\[
M/{\equiv_\lambda}
\cong
\operatorname{Der}(\lambda).
\]
This is the monoid-input Myhill--Nerode theorem without any finite-state
assumption.
\end{example}

\section{Realizations and minimality}
\label{sec:realizations-and-minimality}

Let
\[
k:K\to\mathsf{Beh}_{\mathbb A,\mathbb B}
\]
be a behaviour specification in \(\mathcal E/Z\).

\begin{definition}[Realization]
\label{def:behaviour-realization}
A \textbf{realization} of \(k\) is a Mealy system \(P\), together with a map
\[
i:K\to P
\]
in \(\mathcal E/Z\), such that
\[
\beh_P i=k.
\]
\end{definition}

Assume \(\mathcal E\) is regular. Given a realization
\[
i:K\to P,
\]
define the reachable part of \(P\) from \(K\) as the image in
\(\mathcal E/Z\) of the map
\[
\rho_{P,K}:
A_1\ltimes K
\to
P,
\qquad
(a,k)\mapsto \delta_P(a,i(k)).
\]
Here \(A_1\ltimes K\) is regarded as an object over \(Z\) by the composite
\[
A_1\ltimes K
\xrightarrow{d_K}
\mathsf{Beh}_{\mathbb A,\mathbb B}
\to Z.
\]
The equality
\[
\beh_P\rho_{P,K}=d_K
\]
shows that \(\rho_{P,K}\) is a morphism in \(\mathcal E/Z\). Denote its image
by
\[
\operatorname{Reach}_P(K)\hookrightarrow P.
\]

\begin{proposition}
\label{prop:reachable-part-submealy}
The image
\[
\operatorname{Reach}_P(K)\hookrightarrow P
\]
is a sub-Mealy system of \(P\).
\end{proposition}

\begin{proof}
Let
\[
A_1\ltimes K
\twoheadrightarrow
\operatorname{Reach}_P(K)
\hookrightarrow
P
\]
be the image factorization of \(\rho_{P,K}\). Since regular epimorphisms are
stable under pullback, the induced map
\[
A_1\times_{t_A,A_0,p_\rho}(A_1\ltimes K)
\twoheadrightarrow
A_1\times_{t_A,A_0,p_{\operatorname{Reach}}}
\operatorname{Reach}_P(K)
\]
is a regular epi, where \(p_\rho\) is the \(A_0\)-component of
\(\rho_{P,K}\).

A generalized element of the left-hand side is a triple
\[
(c,a,k)
\]
with
\[
a:u\to v,
\qquad
c:r\to u,
\qquad
p_K(k)=v.
\]
After pulling back along the displayed regular epimorphism, the transition of
\(P\) sends this element to
\[
\delta_P(c,\delta_P(a,i(k))).
\]
By the Mealy transition law,
\[
\delta_P(c,\delta_P(a,i(k)))
=
\delta_P(a\circ c,i(k)).
\]
This lies in the image of \(\rho_{P,K}\), via
\[
(c,a,k)\longmapsto (m_A(c,a),k).
\]
Therefore the transition of \(P\), after pullback along the regular epi,
lands in \(\operatorname{Reach}_P(K)\). Pullback along a regular epimorphism
reflects inclusion of subobjects in a regular category, so the transition of
\(P\) restricts to
\[
A_1\times_{A_0}\operatorname{Reach}_P(K)
\to
\operatorname{Reach}_P(K).
\]

The output map restricts by restricting the output map of \(P\) to
\[
A_1\times_{A_0}\operatorname{Reach}_P(K).
\]
The Mealy equations hold because they hold in \(P\). Hence
\(\operatorname{Reach}_P(K)\) is a sub-Mealy system of \(P\).
\end{proof}

\begin{definition}[Reachable and behaviourally separated realization]
\label{def:reachable-separated-realization}
Let
\[
i:K\to P
\]
be a realization of \(k\).

The realization is \textbf{reachable from \(K\)} if the inclusion
\[
\operatorname{Reach}_P(K)\hookrightarrow P
\]
is an isomorphism.

It is \textbf{behaviourally separated} if the behaviour map
\[
\beh_P:P\to\mathsf{Beh}_{\mathbb A,\mathbb B}
\]
is a monomorphism.
\end{definition}

\begin{remark}[Reachable and separated]
\label{rem:reachable-separated-interpretation}
The two conditions are the internal versions of accessibility and
observability.

Reachability says that every state of \(P\) is obtained by starting from a
specified behaviour in \(K\) and acting by some input arrow of \(\mathbb A\).
Since arrows of \(\mathbb A\) may already encode composite programs, this is
reachability by all available input programs, not necessarily by generator
steps.

Behavioural separation says that the behaviour map
\[
P\to\mathsf{Beh}_{\mathbb A,\mathbb B}
\]
is monic. Thus distinct states of \(P\) have distinct residual-output
behaviours. Equivalently, \(P\) contains no redundant behaviourally equivalent
states.
\end{remark}

\begin{proposition}
\label{prop:reachable-part-maps-to-derivative-closure}
There is a regular epimorphism
\[
\operatorname{Reach}_P(K)
\twoheadrightarrow
\operatorname{Der}(K).
\]
\end{proposition}

\begin{proof}
The behaviour map preserves derivatives. Hence the composite
\[
A_1\ltimes K
\xrightarrow{\rho_{P,K}}
P
\xrightarrow{\beh_P}
\mathsf{Beh}_{\mathbb A,\mathbb B}
\]
agrees with
\[
d_K:
A_1\ltimes K
\to
\mathsf{Beh}_{\mathbb A,\mathbb B}.
\]
Taking images gives a regular epimorphism from the image in \(P\), namely
\[
\operatorname{Reach}_P(K),
\]
onto the image in the final behaviour object, namely
\[
\operatorname{Der}(K).
\]
\end{proof}

\begin{theorem}[Minimal realization]
\label{thm:minimal-realization}
Assume \(\mathcal E\) is regular. If a realization \(P\) of \(K\) is reachable
from \(K\) and behaviourally separated, then
\[
P\cong\operatorname{Der}(K)
\]
as Mealy systems.
\end{theorem}

\begin{proof}
If \(P\) is reachable from \(K\), then
\[
\operatorname{Reach}_P(K)=P.
\]
By Proposition~\ref{prop:reachable-part-maps-to-derivative-closure}, the
behaviour map induces a regular epimorphism
\[
P\twoheadrightarrow\operatorname{Der}(K).
\]
If \(P\) is behaviourally separated, its behaviour map is monic. Since
\[
P\to\operatorname{Der}(K)\hookrightarrow
\mathsf{Beh}_{\mathbb A,\mathbb B}
\]
is the behaviour map of \(P\), the regular epimorphism
\[
P\twoheadrightarrow\operatorname{Der}(K)
\]
is also monic. Therefore it is an isomorphism in \(\mathcal E/Z\).

The map is a Mealy transformation because it is induced by the behaviour map,
and the inverse is a Mealy transformation because the Mealy structure on
\(\operatorname{Der}(K)\) is the restricted structure inherited from the final
behaviour system. Hence the isomorphism is an isomorphism of Mealy systems.
\end{proof}

\begin{remark}[Internal Myhill--Nerode principle]
\label{rem:internal-myhill-nerode-principle}
The derivative closure
\[
\operatorname{Der}(K)
\]
is the canonical reachable separated realization. Any other reachable
realization maps onto it by the behaviour map. If the other realization is
also separated, that map is forced to be an isomorphism.

Thus the Myhill--Nerode object is simultaneously:
\begin{enumerate}
\item the reachable residual subobject of the final behaviour object;
\item the quotient of residual descriptions by Nerode equivalence;
\item the unique reachable and separated realization of the specification.
\end{enumerate}
This is the internal-categorical version of the usual minimal automaton
theorem.
\end{remark}

\begin{example}[Minimal realization in the monoid case]
\label{ex:monoid-minimal-realization}
For a behaviour
\[
\lambda\in\mathsf{Beh}_{M,N},
\]
the minimal realization has carrier
\[
\operatorname{Der}(\lambda)
=
\{\partial_r\lambda\mid r\in M\}.
\]
Its action and output are
\[
\partial_r\lambda\cdot m=\partial_{rm}\lambda,
\qquad
\omega(m,\partial_r\lambda)=\lambda(r,m).
\]
Every reachable and behaviourally separated realization of \(\lambda\) is
isomorphic to this residual-state system.
\end{example}

\section{The finite-state corollary}
\label{sec:finite-state-corollary}

The general Nerode realization theorem does not require finiteness. To recover
the finite-state form, choose a class
\[
\mathsf{Fin}_Z\subseteq\mathcal E/Z
\]
of finite-state carrier objects. Assume that \(\mathsf{Fin}_Z\) is closed
under isomorphism, subobjects, regular images, and effective quotients of
internal equivalence relations.

\begin{definition}[Finite-state realization]
\label{def:finite-state-realization}
A realization \(P\) of a specification \(K\) is \textbf{finite-state} if its
carrier lies in
\[
\mathsf{Fin}_Z.
\]
\end{definition}

\begin{corollary}[Finite-state Myhill--Nerode]
\label{cor:finite-state-myhill-nerode}
A specification
\[
k:K\to\mathsf{Beh}_{\mathbb A,\mathbb B}
\]
has a finite-state realization if and only if
\[
\operatorname{Der}(K)\in\mathsf{Fin}_Z.
\]
Equivalently, when \(\mathcal E\) is Barr-exact, it has a finite-state
realization if and only if its Nerode quotient
\[
(A_1\ltimes K)/{\equiv_K}
\]
lies in
\[
\mathsf{Fin}_Z.
\]
\end{corollary}

\begin{proof}
If \(K\) has a finite-state realization \(P\), then its reachable part
\[
\operatorname{Reach}_P(K)
\]
is a subobject of \(P\), hence lies in \(\mathsf{Fin}_Z\). The derivative
closure
\[
\operatorname{Der}(K)
\]
is a regular image of \(\operatorname{Reach}_P(K)\), so it also lies in
\(\mathsf{Fin}_Z\).

Conversely, if
\[
\operatorname{Der}(K)\in\mathsf{Fin}_Z,
\]
then \(\operatorname{Der}(K)\) itself is a finite-state realization of \(K\).
The equivalence with the Nerode quotient follows from
Theorem~\ref{thm:general-nerode-realization} and closure of \(\mathsf{Fin}_Z\)
under isomorphism.
\end{proof}

\begin{remark}[Finiteness is external to the construction]
\label{rem:finiteness-external-to-construction}
The construction of
\[
\operatorname{Der}(K)
\]
does not require a finiteness hypothesis. Finiteness enters only after
choosing a class
\[
\mathsf{Fin}_Z
\]
of carriers that one wants to regard as finite-state. Closure of
\(\mathsf{Fin}_Z\) under subobjects and regular images gives the
derivative-closure criterion. Closure under effective quotients is used for
the equivalent Nerode-quotient formulation. The finite-state Myhill--Nerode
statement says that \(K\) is finitely realizable exactly when its canonical
residual realization is finite.
\end{remark}

\begin{example}[Finite-state monoid case]
\label{ex:finite-state-monoid-case}
For a behaviour
\[
\lambda\in\mathsf{Beh}_{M,N},
\]
one has
\[
\lambda
\text{ is finite-state realizable}
\]
if and only if
\[
\operatorname{Der}(\lambda)
=
\{\partial_r\lambda\mid r\in M\}
\]
is finite. Equivalently,
\[
M/{\equiv_\lambda}
\]
is finite.
\end{example}

\section{Comparison with the raw Set-coalgebra functor}
\label{sec:comparison-raw-set-mealy-functor}

In ordinary coalgebraic automata theory, a deterministic Mealy machine with
input set \(I\) and output set \(O\) is often represented as a coalgebra for
the functor
\[
X\mapsto (O\times X)^I.
\]

The one-object specialization of the present theory should not be confused
with this raw Set-based functor. Here the input object is not a bare set. It
is a monoid \(M\). Likewise the output object is a monoid \(N\). The
corresponding systems are therefore not arbitrary maps
\[
M\times P\to P\times N.
\]
They are precisely those maps satisfying the action and cocycle equations:
\[
x\cdot(mn)=(x\cdot m)\cdot n,
\]
\[
\omega(mn,x)=\omega(m,x)\omega(n,x\cdot m).
\]

Thus the final behaviour object is not the set of arbitrary \(M\)-branching
\(N\)-labelled trees. It is the subobject of coherent residual-output
behaviours:
\[
\lambda:M\times M\to N
\]
satisfying
\[
\lambda(r,1)=1,
\qquad
\lambda(r,mn)=\lambda(r,m)\lambda(rm,n).
\]

When \(M\) is specialized to a free monoid \(I^\ast\), this coherent
monoid-input behaviour object has the generator-level normal form
\[
N^{I^+}.
\]
This is the classical Mealy final-behaviour object. Thus the present theory
does not replace the classical object; rather, it refines it by first treating
inputs as an arbitrary monoid and then recovering the classical exponential
form in the free-input specialization.

\begin{remark}
\label{rem:raw-versus-monoid-input}
If \(M\) is itself a free monoid of input words and \(N\) is a monoid of
outputs, then the present theory describes word-level Mealy behaviour whose
outputs compose by the cocycle law. Restricting to generator inputs recovers
the ordinary letter-step Mealy machine. This distinction is important: the
native theory is monoid-input, while the classical functor
\[
X\mapsto (N\times X)^I
\]
is generator-input.
\end{remark}

\section{Main consolidation}
\label{sec:mealy-comonad-main-consolidation}

\begin{theorem}[Mealy comonad and Nerode realization]
\label{thm:mealy-comonad-main-consolidation}
Let
\[
\mathcal E
\]
be locally presentable and locally cartesian closed, and let
\[
\mathbb A,\mathbb B
\]
be internal categories in \(\mathcal E\). Then:

\begin{enumerate}
\item The category
\[
\mathsf{Mealy}(\mathbb A,\mathbb B)
\]
of Mealy morphisms and Mealy transformations is locally presentable.

\item The forgetful functor
\[
U_{\mathbb A,\mathbb B}:
\mathsf{Mealy}(\mathbb A,\mathbb B)
\to
\mathcal E/(A_0\times B_0)
\]
has a right adjoint
\[
R_{\mathbb A,\mathbb B}.
\]

\item The composite
\[
\mathbb M_{\mathbb A,\mathbb B}
=
U_{\mathbb A,\mathbb B}R_{\mathbb A,\mathbb B}
\]
is a comonad on
\[
\mathcal E/(A_0\times B_0).
\]

\item There is an equivalence
\[
\mathsf{Mealy}(\mathbb A,\mathbb B)
\simeq
\mathbf{Coalg}(\mathbb M_{\mathbb A,\mathbb B}).
\]

\item The final Mealy behaviour system is
\[
\mathsf{Beh}_{\mathbb A,\mathbb B}
=
R_{\mathbb A,\mathbb B}(1_{A_0\times B_0}).
\]

\item Every Mealy system \(P\) has a unique behaviour map
\[
\beh_P:P\to\mathsf{Beh}_{\mathbb A,\mathbb B}.
\]

\item Behaviours admit a residual-slice normal form. For a generalized state
\(x\) of a Mealy system, with
\[
p(x)=v,
\qquad
q(x)=y,
\]
the behaviour of \(x\) is represented by the pointed functor
\[
H_x:\mathbb A/v\to\mathbb B
\]
defined by
\[
H_x(r)=q(\delta(r,x)),
\]
and
\[
H_x(a:r\circ a\to r)=\omega(a,\delta(r,x)).
\]
The value at the identity is
\[
H_x(1_v)=y.
\]

\item Derivatives are residual precomposition. If
\[
r:u\to v,
\]
then
\[
\partial_rH=H\circ r_\ast,
\]
where
\[
r_\ast:\mathbb A/u\to\mathbb A/v
\]
is postcomposition by \(r\).

\item In the external one-object Set case, where
\[
\mathcal E=\mathbf{Set},
\qquad
\mathbb A=M,
\qquad
\mathbb B=N
\]
are monoids, the final behaviour object is
\[
\mathsf{Beh}_{M,N}
\cong
[M/\ast,N],
\]
equivalently the set of residual cocycles
\[
\lambda:M\times M\to N
\]
satisfying
\[
\lambda(r,1)=1,
\qquad
\lambda(r,mn)=\lambda(r,m)\lambda(rm,n).
\]

\item If moreover
\[
M=I^\ast
\]
is the free monoid on a set \(I\), then
\[
\mathsf{Beh}_{I^\ast,N}
\cong
N^{I^+},
\]
recovering the standard final behaviour object for classical deterministic
Mealy machines.

\item If \(\mathcal E\) is regular, the image of \(\beh_P\) is the realized
behaviour object
\[
\operatorname{Real}(P).
\]

\item If \(\mathcal E\) is regular, then for every behaviour specification
\[
K\to\mathsf{Beh}_{\mathbb A,\mathbb B},
\]
the derivative closure
\[
\operatorname{Der}(K)
\]
is the smallest sub-Mealy system of the final behaviour object containing the
image of \(K\).

\item If \(\mathcal E\) is Barr-exact, then for every behaviour specification
\[
K\to\mathsf{Beh}_{\mathbb A,\mathbb B},
\]
the Nerode quotient is canonically isomorphic to the derivative closure:
\[
(A_1\ltimes K)/{\equiv_K}
\cong
\operatorname{Der}(K).
\]
With the Mealy structure transported across this isomorphism, the Nerode
quotient and the derivative closure are the same Mealy realization.

\item If \(\mathcal E\) is regular, then a reachable and behaviourally
separated realization of \(K\) is canonically isomorphic, as a Mealy system,
to
\[
\operatorname{Der}(K).
\]

\item After choosing a finite-state class
\[
\mathsf{Fin}_{A_0\times B_0}\subseteq\mathcal E/(A_0\times B_0)
\]
closed under isomorphism, subobjects, regular images, and effective quotients,
the finite-state Myhill--Nerode theorem follows as a corollary:
\[
K \text{ has a finite-state realization}
\quad\Longleftrightarrow\quad
\operatorname{Der}(K)
\text{ is finite-state}.
\]
When \(\mathcal E\) is Barr-exact, this is equivalently the finiteness of the
Nerode quotient.
\end{enumerate}
\end{theorem}

\begin{proof}
Part (1) is Proposition~\ref{prop:mealy-category-locally-presentable}. Part
(2) is Theorem~\ref{thm:cofree-mealy-systems}. Part (3) is
Definition~\ref{def:mealy-comonad}. Part (4) is
Theorem~\ref{thm:mealy-systems-are-comonad-coalgebras}. Part (5) is
Definition~\ref{def:behaviour-object} and
Theorem~\ref{thm:final-behaviour-coalgebra}. Part (6) is
Definition~\ref{def:behaviour-map}. Part (7) is
Proposition~\ref{prop:behaviours-as-residual-slice-functors}. Part (8) is
Proposition~\ref{prop:derivatives-as-residual-precomposition}. Part (9) is
Proposition~\ref{prop:one-object-residual-slice-form} and
Corollary~\ref{cor:final-monoid-behaviour}. Part (10) is
Corollary~\ref{cor:free-input-classical-behaviour-object}. Part (11) is
Definition~\ref{def:realized-behaviour-object} and
Proposition~\ref{prop:realized-behaviour-submealy}. Part (12) is
Theorem~\ref{thm:derivative-closure}. Part (13) is
Theorem~\ref{thm:general-nerode-realization}. Part (14) is
Theorem~\ref{thm:minimal-realization}. Part (15) is
Corollary~\ref{cor:finite-state-myhill-nerode}.
\end{proof}

\begin{remark}[Final one-object Set summary]
\label{rem:final-one-object-summary}
In the external one-object Set case with
\[
\mathcal E=\mathbf{Set},
\qquad
\mathbb A=M,
\qquad
\mathbb B=N,
\]
the chapter specializes to the following data.

A system is a right \(M\)-set with an \(N\)-valued cocycle. The underlying
carrier of the Mealy comonad on sets is
\[
\mathbb M_{M,N}(X)
=
\left\{
(\ell,\lambda)
\ \middle|\
\ell:M\to X,\
\lambda:M\times M\to N,\
\lambda(r,1)=1,\
\lambda(r,mn)=\lambda(r,m)\lambda(rm,n)
\right\}.
\]
The final behaviour object is obtained by taking \(X=1\):
\[
\mathsf{Beh}_{M,N}
=
\left\{
\lambda:M\times M\to N
\ \middle|\
\lambda(r,1)=1,\
\lambda(r,mn)=\lambda(r,m)\lambda(rm,n)
\right\}.
\]
Derivatives are residual shifts:
\[
(\partial_r\lambda)(s,m)=\lambda(rs,m).
\]
The Nerode equivalence of a behaviour \(\lambda\) is
\[
r\equiv_\lambda s
\quad\Longleftrightarrow\quad
\partial_r\lambda=\partial_s\lambda,
\]
and the general Nerode theorem says
\[
M/{\equiv_\lambda}
\cong
\operatorname{Der}(\lambda).
\]

If \(M=I^\ast\) is free, then the same behaviour object has the classical
normal form
\[
\mathsf{Beh}_{I^\ast,N}
\cong
N^{I^+}.
\]
Thus the coherent monoid-input final behaviour object recovers the ordinary
Mealy final behaviour object in the free-input specialization.
\end{remark}